\chapter{Logistische Regression}

\section{Ziel dieses Kapitels}

In den vorherigen Kapiteln haben wir lineare Regression (linear regression), Maximum-Likelihood-Schätzung (maximum likelihood estimation), Maximum-a-posteriori-Schätzung (maximum a-posteriori estimation) und negative Log-Likelihood kennengelernt.
Lineare Regression eignet sich für kontinuierliche Zielwerte.
In diesem Kapitel betrachten wir nun Klassifikation (classification), also Probleme mit diskreten Zielwerten.

Die logistische Regression (logistic regression) ist eines der wichtigsten Grundmodelle für binäre Klassifikation (binary classification).
Trotz des Namens ist sie ein Klassifikationsmodell und kein Regressionsmodell im üblichen Sinn.

\begin{conceptbox}
Logistische Regression (logistic regression) ist ein probabilistisches Klassifikationsmodell.

Für eine Eingabe \(\vec{x}\) berechnet das Modell eine Wahrscheinlichkeit

\[
    p(t=1 \mid \vec{x}).
\]

Diese Wahrscheinlichkeit wird dann verwendet, um zwischen zwei Klassen zu entscheiden.
\end{conceptbox}

Nach diesem Kapitel solltest du folgende Fragen beantworten können:

\begin{itemize}
    \item Was ist binäre Klassifikation (binary classification)?
    \item Warum ist lineare Regression für Klassifikation ungeeignet?
    \item Was ist die Sigmoid-Funktion (sigmoid function)?
    \item Was ist der Zusammenhang zwischen Sigmoid, Odds und Logit?
    \item Wie modelliert logistische Regression \(p(t=1\mid\vec{x})\)?
    \item Wie entsteht die Bernoulli-Likelihood?
    \item Warum führt Maximum-Likelihood zu Kreuzentropie (cross-entropy)?
    \item Wie leitet man den Gradienten der logistischen Regression her?
    \item Warum gibt es keine geschlossene Lösung wie bei linearer Regression?
    \item Wie trainiert man logistische Regression mit Gradientenabstieg (gradient descent)?
    \item Was ist eine Entscheidungsgrenze (decision boundary)?
    \item Wie erweitert man logistische Regression auf mehrere Klassen?
    \item Wie hängt Regularisierung mit MAP zusammen?
\end{itemize}

\section{Von Regression zu Klassifikation}

Bei Regression (regression) ist der Zielwert kontinuierlich:

\[
    t \in \mathbb{R}.
\]

Bei Klassifikation (classification) ist der Zielwert diskret.
Bei binärer Klassifikation gibt es genau zwei Klassen:

\[
    t \in \{0,1\}.
\]

\begin{definitionbox}
Binäre Klassifikation (binary classification) ist ein überwachtes Lernproblem, bei dem jeder Datenpunkt zu einer von zwei Klassen gehört.

Der Trainingsdatensatz hat die Form

\[
    \mathcal{D}
    =
    \{(\vec{x}_1,t_1),\dots,(\vec{x}_N,t_N)\},
\]

wobei

\[
    \vec{x}_n \in \mathbb{R}^D
\]

und

\[
    t_n \in \{0,1\}.
\]
\end{definitionbox}

Beispiele für binäre Klassifikation:

\begin{itemize}
    \item E-Mail: Spam oder Nicht-Spam,
    \item medizinischer Test: krank oder gesund,
    \item Bild: Katze oder keine Katze,
    \item Kreditentscheidung: Ausfall oder kein Ausfall,
    \item Wohnungspreis-Beispiel: teuer oder nicht teuer.
\end{itemize}

\section{Warum nicht einfach lineare Regression?}

Man könnte versuchen, ein lineares Regressionsmodell zu verwenden:

\[
    y(\vec{x},\vec{w})
    =
    \vec{w}^{\top}\vec{x}.
\]

Dann könnte man sagen:

\[
    y(\vec{x},\vec{w}) \geq 0{,}5
    \quad
    \Rightarrow
    \quad
    \hat{t}=1,
\]

und

\[
    y(\vec{x},\vec{w}) < 0{,}5
    \quad
    \Rightarrow
    \quad
    \hat{t}=0.
\]

Das ist aber problematisch.

\begin{conceptbox}
Lineare Regression ist für Klassifikation ungeeignet, weil die Ausgabe

\[
    y(\vec{x},\vec{w})
    =
    \vec{w}^{\top}\vec{x}
\]

beliebige reelle Werte annehmen kann.

Für eine Wahrscheinlichkeit brauchen wir aber Werte im Intervall

\[
    [0,1].
\]
\end{conceptbox}

Eine Wahrscheinlichkeit muss erfüllen:

\[
    0 \leq p(t=1\mid\vec{x}) \leq 1.
\]

Ein lineares Modell garantiert das nicht.
Zum Beispiel könnte gelten:

\[
    \vec{w}^{\top}\vec{x}
    =
    -3
\]

oder

\[
    \vec{w}^{\top}\vec{x}
    =
    4{,}7.
\]

Beides kann keine Wahrscheinlichkeit sein.

\section{Die Grundidee der logistischen Regression}

Die Idee der logistischen Regression ist:

\begin{enumerate}
    \item Berechne zuerst einen linearen Score:
    \[
        a
        =
        \vec{w}^{\top}\vec{x}.
    \]

    \item Wandle diesen Score mit einer Sigmoid-Funktion in eine Wahrscheinlichkeit um:
    \[
        y
        =
        \sigma(a).
    \]
\end{enumerate}

\begin{definitionbox}
Das Modell der binären logistischen Regression lautet:

\[
    y(\vec{x},\vec{w})
    =
    \sigma(\vec{w}^{\top}\vec{x}),
\]

wobei

\[
    y(\vec{x},\vec{w})
    =
    p(t=1\mid\vec{x},\vec{w}).
\]
\end{definitionbox}

Der Wert

\[
    a
    =
    \vec{w}^{\top}\vec{x}
\]

heißt Aktivierung (activation), Score oder Logit.

Die Ausgabe

\[
    y
    =
    \sigma(a)
\]

ist eine Wahrscheinlichkeit.

\section{Die Sigmoid-Funktion}

Die Sigmoid-Funktion (sigmoid function) ist definiert als

\[
    \sigma(a)
    =
    \frac{1}{1+\exp(-a)}.
\]

\begin{definitionbox}
Die Sigmoid-Funktion ist

\[
    \sigma(a)
    =
    \frac{1}{1+\exp(-a)}.
\]

Sie bildet reelle Zahlen auf das Intervall \((0,1)\) ab:

\[
    \sigma:\mathbb{R}\to(0,1).
\]
\end{definitionbox}

Wichtige Eigenschaften:

\[
    \lim_{a\to\infty}\sigma(a)=1,
\]

\[
    \lim_{a\to-\infty}\sigma(a)=0,
\]

und

\[
    \sigma(0)=\frac{1}{2}.
\]

\begin{conceptbox}
Die Sigmoid-Funktion macht aus einem beliebigen linearen Score \(a\in\mathbb{R}\) eine gültige Wahrscheinlichkeit.

\[
    a \in \mathbb{R}
    \quad
    \longrightarrow
    \quad
    \sigma(a)\in(0,1).
\]
\end{conceptbox}

\section{Interpretation der Sigmoid-Ausgabe}

Bei logistischer Regression setzen wir:

\[
    p(t=1\mid\vec{x},\vec{w})
    =
    y(\vec{x},\vec{w})
    =
    \sigma(\vec{w}^{\top}\vec{x}).
\]

Da es nur zwei Klassen gibt, muss gelten:

\[
    p(t=0\mid\vec{x},\vec{w})
    =
    1-p(t=1\mid\vec{x},\vec{w}).
\]

Also:

\[
    p(t=0\mid\vec{x},\vec{w})
    =
    1-y(\vec{x},\vec{w}).
\]

\begin{definitionbox}
Für binäre logistische Regression gilt:

\[
    p(t=1\mid\vec{x},\vec{w})
    =
    y,
\]

und

\[
    p(t=0\mid\vec{x},\vec{w})
    =
    1-y,
\]

wobei

\[
    y
    =
    \sigma(\vec{w}^{\top}\vec{x}).
\]
\end{definitionbox}

\section{Entscheidungsregel}

Aus der Wahrscheinlichkeit kann eine Klassenentscheidung gemacht werden.
Die Standardregel ist:

\[
    \hat{t}
    =
    \begin{cases}
        1, & y \geq 0{,}5, \\
        0, & y < 0{,}5.
    \end{cases}
\]

Da

\[
    \sigma(0)=0{,}5
\]

und die Sigmoid-Funktion monoton wachsend ist, gilt:

\[
    y \geq 0{,}5
    \quad
    \Longleftrightarrow
    \quad
    \vec{w}^{\top}\vec{x} \geq 0.
\]

\begin{conceptbox}
Die Entscheidungsgrenze (decision boundary) der logistischen Regression ist

\[
    \vec{w}^{\top}\vec{x}=0.
\]

Auf der einen Seite entscheidet das Modell Klasse \(1\), auf der anderen Seite Klasse \(0\).
\end{conceptbox}

\section{Bias-Term}

Wie bei linearer Regression verwenden wir meistens einen Bias-Term.
Für eine Eingabe

\[
    \vec{x}
    =
    \begin{pmatrix}
        x_1 \\
        \vdots \\
        x_D
    \end{pmatrix}
\]

erweitern wir:

\[
    \tilde{\vec{x}}
    =
    \begin{pmatrix}
        1 \\
        x_1 \\
        \vdots \\
        x_D
    \end{pmatrix}.
\]

Der Gewichtsvektor ist:

\[
    \vec{w}
    =
    \begin{pmatrix}
        w_0 \\
        w_1 \\
        \vdots \\
        w_D
    \end{pmatrix}.
\]

Dann gilt:

\[
    a
    =
    \vec{w}^{\top}\tilde{\vec{x}}
    =
    w_0+\sum_{j=1}^{D}w_jx_j.
\]

\begin{notebox}
Der Bias-Term \(w_0\) verschiebt die Entscheidungsgrenze.

Ohne Bias müsste die Entscheidungsgrenze durch den Ursprung gehen.
\end{notebox}

Im Folgenden schreiben wir oft einfach \(\vec{x}\) statt \(\tilde{\vec{x}}\), wenn der Bias-Eintrag bereits enthalten ist.

\section{Odds und Logit}

Die logistische Regression hat eine wichtige Interpretation über Odds und Logit.

Sei

\[
    y
    =
    p(t=1\mid\vec{x}).
\]

Dann ist

\[
    1-y
    =
    p(t=0\mid\vec{x}).
\]

Die Odds sind definiert als:

\[
    \frac{y}{1-y}.
\]

\begin{definitionbox}
Die Odds (odds) sind das Verhältnis der Wahrscheinlichkeit für Klasse \(1\) zur Wahrscheinlichkeit für Klasse \(0\):

\[
    \mathrm{odds}
    =
    \frac{p(t=1\mid\vec{x})}{p(t=0\mid\vec{x})}
    =
    \frac{y}{1-y}.
\]
\end{definitionbox}

Der Logit ist der Logarithmus der Odds:

\[
    \operatorname{logit}(y)
    =
    \log
    \left(
        \frac{y}{1-y}
    \right).
\]

Bei logistischer Regression gilt:

\[
    y
    =
    \sigma(a)
    =
    \frac{1}{1+\exp(-a)}.
\]

Dann ist:

\[
    1-y
    =
    1-
    \frac{1}{1+\exp(-a)}
    =
    \frac{\exp(-a)}{1+\exp(-a)}.
\]

Also:

\[
    \frac{y}{1-y}
    =
    \frac{
        \frac{1}{1+\exp(-a)}
    }{
        \frac{\exp(-a)}{1+\exp(-a)}
    }
    =
    \frac{1}{\exp(-a)}
    =
    \exp(a).
\]

Damit:

\[
    \log
    \left(
        \frac{y}{1-y}
    \right)
    =
    \log(\exp(a))
    =
    a.
\]

Da \(a=\vec{w}^{\top}\vec{x}\), folgt:

\[
    \operatorname{logit}(y)
    =
    \vec{w}^{\top}\vec{x}.
\]

\begin{conceptbox}
Logistische Regression modelliert den Logit linear:

\[
    \log
    \left(
        \frac{p(t=1\mid\vec{x})}{p(t=0\mid\vec{x})}
    \right)
    =
    \vec{w}^{\top}\vec{x}.
\]

Die Wahrscheinlichkeit selbst ist nicht linear in \(\vec{x}\), aber die Log-Odds sind linear.
\end{conceptbox}

\section{Bernoulli-Modell für die Zielwerte}

Da

\[
    t_n \in \{0,1\}
\]

ist, modellieren wir \(t_n\) mit einer Bernoulli-Verteilung (Bernoulli distribution).

Für einen Datenpunkt gilt:

\[
    p(t_n\mid\vec{x}_n,\vec{w})
    =
    y_n^{t_n}
    (1-y_n)^{1-t_n},
\]

wobei

\[
    y_n
    =
    y(\vec{x}_n,\vec{w})
    =
    \sigma(\vec{w}^{\top}\vec{x}_n).
\]

\begin{definitionbox}
Das probabilistische Modell der logistischen Regression ist:

\[
    p(t\mid\vec{x},\vec{w})
    =
    y^t(1-y)^{1-t},
\]

mit

\[
    y
    =
    \sigma(\vec{w}^{\top}\vec{x}).
\]
\end{definitionbox}

Prüfen wir die beiden Fälle:

Falls

\[
    t=1,
\]

dann:

\[
    p(t=1\mid\vec{x},\vec{w})
    =
    y^1(1-y)^0
    =
    y.
\]

Falls

\[
    t=0,
\]

dann:

\[
    p(t=0\mid\vec{x},\vec{w})
    =
    y^0(1-y)^1
    =
    1-y.
\]

\section{Likelihood der logistischen Regression}

Für einen Datensatz

\[
    \mathcal{D}
    =
    \{(\vec{x}_1,t_1),\dots,(\vec{x}_N,t_N)\}
\]

nehmen wir i.i.d. Daten an.
Dann ist die Likelihood:

\[
    p(\vec{t}\mid X,\vec{w})
    =
    \prod_{n=1}^{N}
    p(t_n\mid\vec{x}_n,\vec{w}).
\]

Einsetzen des Bernoulli-Modells ergibt:

\[
    p(\vec{t}\mid X,\vec{w})
    =
    \prod_{n=1}^{N}
    y_n^{t_n}
    (1-y_n)^{1-t_n}.
\]

\begin{definitionbox}
Die Likelihood der logistischen Regression lautet:

\[
    L(\vec{w})
    =
    \prod_{n=1}^{N}
    y_n^{t_n}
    (1-y_n)^{1-t_n},
\]

wobei

\[
    y_n
    =
    \sigma(\vec{w}^{\top}\vec{x}_n).
\]
\end{definitionbox}

\section{Log-Likelihood}

Die Log-Likelihood ist:

\[
    \ell(\vec{w})
    =
    \log L(\vec{w}).
\]

Also:

\[
    \ell(\vec{w})
    =
    \log
    \prod_{n=1}^{N}
    y_n^{t_n}
    (1-y_n)^{1-t_n}.
\]

Der Logarithmus macht aus dem Produkt eine Summe:

\[
    \ell(\vec{w})
    =
    \sum_{n=1}^{N}
    \log
    \left[
        y_n^{t_n}
        (1-y_n)^{1-t_n}
    \right].
\]

Mit Logarithmusregeln:

\[
    \ell(\vec{w})
    =
    \sum_{n=1}^{N}
    \left[
        t_n\log y_n
        +
        (1-t_n)\log(1-y_n)
    \right].
\]

\begin{definitionbox}
Die Log-Likelihood der logistischen Regression ist:

\[
    \ell(\vec{w})
    =
    \sum_{n=1}^{N}
    \left[
        t_n\log y_n
        +
        (1-t_n)\log(1-y_n)
    \right].
\]
\end{definitionbox}

\section{Negative Log-Likelihood und Kreuzentropie}

In maschinellem Lernen minimieren wir meistens eine Verlustfunktion.
Daher betrachten wir die negative Log-Likelihood:

\[
    E(\vec{w})
    =
    -\ell(\vec{w}).
\]

Damit:

\[
    E(\vec{w})
    =
    -\sum_{n=1}^{N}
    \left[
        t_n\log y_n
        +
        (1-t_n)\log(1-y_n)
    \right].
\]

\begin{definitionbox}
Die Fehlerfunktion der logistischen Regression ist die binäre Kreuzentropie (binary cross-entropy):

\[
    E(\vec{w})
    =
    -\sum_{n=1}^{N}
    \left[
        t_n\log y_n
        +
        (1-t_n)\log(1-y_n)
    \right].
\]
\end{definitionbox}

Für einen einzelnen Datenpunkt ist der Verlust:

\[
    E_n(\vec{w})
    =
    -\left[
        t_n\log y_n
        +
        (1-t_n)\log(1-y_n)
    \right].
\]

\begin{conceptbox}
Kreuzentropie (cross-entropy) ist die negative Log-Likelihood eines Bernoulli-Modells.

Sie ist also keine willkürlich gewählte Fehlerfunktion, sondern folgt aus Maximum-Likelihood.
\end{conceptbox}

\section{Intuition der Kreuzentropie}

Betrachten wir einen einzelnen Datenpunkt.

Falls

\[
    t=1,
\]

dann ist der Verlust:

\[
    E
    =
    -\log y.
\]

Wenn das Modell

\[
    y\approx 1
\]

vorhersagt, dann ist

\[
    -\log y \approx 0.
\]

Wenn das Modell aber

\[
    y\approx 0
\]

vorhersagt, dann wird

\[
    -\log y
\]

sehr groß.

Falls

\[
    t=0,
\]

dann ist der Verlust:

\[
    E
    =
    -\log(1-y).
\]

Wenn das Modell

\[
    y\approx 0
\]

vorhersagt, ist der Verlust klein.
Wenn das Modell

\[
    y\approx 1
\]

vorhersagt, ist der Verlust groß.

\begin{conceptbox}
Kreuzentropie bestraft sichere falsche Vorhersagen besonders stark.

Das ist sinnvoll, weil ein Modell nicht nur falsch liegt, sondern auch sehr sicher falsch sein kann.
\end{conceptbox}

\section{Ableitung der Sigmoid-Funktion}

Für den Gradienten der logistischen Regression brauchen wir die Ableitung der Sigmoid-Funktion.

Die Sigmoid-Funktion ist:

\[
    \sigma(a)
    =
    \frac{1}{1+\exp(-a)}.
\]

Wir schreiben:

\[
    \sigma(a)
    =
    (1+\exp(-a))^{-1}.
\]

Ableiten:

\[
    \frac{d\sigma}{da}
    =
    -(1+\exp(-a))^{-2}
    \cdot
    (-\exp(-a)).
\]

Also:

\[
    \frac{d\sigma}{da}
    =
    \frac{\exp(-a)}{(1+\exp(-a))^2}.
\]

Nun gilt:

\[
    \sigma(a)
    =
    \frac{1}{1+\exp(-a)},
\]

und

\[
    1-\sigma(a)
    =
    \frac{\exp(-a)}{1+\exp(-a)}.
\]

Daher:

\[
    \sigma(a)(1-\sigma(a))
    =
    \frac{1}{1+\exp(-a)}
    \cdot
    \frac{\exp(-a)}{1+\exp(-a)}
    =
    \frac{\exp(-a)}{(1+\exp(-a))^2}.
\]

Somit:

\[
    \frac{d\sigma}{da}
    =
    \sigma(a)(1-\sigma(a)).
\]

\begin{definitionbox}
Die Ableitung der Sigmoid-Funktion ist:

\[
    \sigma'(a)
    =
    \sigma(a)(1-\sigma(a)).
\]
\end{definitionbox}

Wenn

\[
    y=\sigma(a),
\]

dann gilt:

\[
    \frac{dy}{da}
    =
    y(1-y).
\]

\section{Gradient für einen Datenpunkt}

Für einen einzelnen Datenpunkt ist der Verlust:

\[
    E_n(\vec{w})
    =
    -\left[
        t_n\log y_n
        +
        (1-t_n)\log(1-y_n)
    \right],
\]

wobei

\[
    y_n
    =
    \sigma(a_n)
\]

und

\[
    a_n
    =
    \vec{w}^{\top}\vec{x}_n.
\]

Wir möchten den Gradienten nach \(\vec{w}\) berechnen:

\[
    \nabla_{\vec{w}}E_n(\vec{w}).
\]

Zuerst leiten wir nach \(y_n\) ab:

\[
    \frac{\partial E_n}{\partial y_n}
    =
    -\left[
        \frac{t_n}{y_n}
        -
        \frac{1-t_n}{1-y_n}
    \right].
\]

Also:

\[
    \frac{\partial E_n}{\partial y_n}
    =
    -\frac{t_n}{y_n}
    +
    \frac{1-t_n}{1-y_n}.
\]

Auf gemeinsamen Nenner gebracht:

\[
    \frac{\partial E_n}{\partial y_n}
    =
    \frac{
        -t_n(1-y_n)
        +
        (1-t_n)y_n
    }{
        y_n(1-y_n)
    }.
\]

Der Zähler ist:

\[
    -t_n+t_ny_n+y_n-t_ny_n
    =
    y_n-t_n.
\]

Also:

\[
    \frac{\partial E_n}{\partial y_n}
    =
    \frac{y_n-t_n}{y_n(1-y_n)}.
\]

Nun:

\[
    \frac{\partial y_n}{\partial a_n}
    =
    y_n(1-y_n).
\]

Mit der Kettenregel:

\[
    \frac{\partial E_n}{\partial a_n}
    =
    \frac{\partial E_n}{\partial y_n}
    \frac{\partial y_n}{\partial a_n}.
\]

Einsetzen:

\[
    \frac{\partial E_n}{\partial a_n}
    =
    \frac{y_n-t_n}{y_n(1-y_n)}
    y_n(1-y_n).
\]

Damit:

\[
    \frac{\partial E_n}{\partial a_n}
    =
    y_n-t_n.
\]

Schließlich gilt:

\[
    a_n
    =
    \vec{w}^{\top}\vec{x}_n.
\]

Also:

\[
    \frac{\partial a_n}{\partial \vec{w}}
    =
    \vec{x}_n.
\]

Damit:

\[
    \nabla_{\vec{w}}E_n(\vec{w})
    =
    (y_n-t_n)\vec{x}_n.
\]

\begin{definitionbox}
Für einen einzelnen Datenpunkt ist der Gradient der Kreuzentropie:

\[
    \nabla_{\vec{w}}E_n(\vec{w})
    =
    (y_n-t_n)\vec{x}_n.
\]

Dabei ist

\[
    y_n
    =
    \sigma(\vec{w}^{\top}\vec{x}_n).
\]
\end{definitionbox}

\begin{conceptbox}
Die Kombination aus Sigmoid und Kreuzentropie führt zu einem besonders einfachen Gradienten:

\[
    \text{Gradient}
    =
    \text{Vorhersagefehler}
    \cdot
    \text{Eingabe}.
\]

Also:

\[
    (y_n-t_n)\vec{x}_n.
\]
\end{conceptbox}

\section{Gradient für den gesamten Datensatz}

Die gesamte Fehlerfunktion ist:

\[
    E(\vec{w})
    =
    \sum_{n=1}^{N}
    E_n(\vec{w}).
\]

Daher ist der Gradient:

\[
    \nabla_{\vec{w}}E(\vec{w})
    =
    \sum_{n=1}^{N}
    \nabla_{\vec{w}}E_n(\vec{w}).
\]

Mit dem Einzeldatenpunkt-Gradienten:

\[
    \nabla_{\vec{w}}E(\vec{w})
    =
    \sum_{n=1}^{N}
    (y_n-t_n)\vec{x}_n.
\]

\begin{definitionbox}
Der Gradient der logistischen Regression ist:

\[
    \nabla_{\vec{w}}E(\vec{w})
    =
    \sum_{n=1}^{N}
    (y_n-t_n)\vec{x}_n.
\]
\end{definitionbox}

\section{Matrixform des Gradienten}

Wir definieren die Designmatrix:

\[
    X
    =
    \begin{pmatrix}
        \vec{x}_1^{\top} \\
        \vec{x}_2^{\top} \\
        \vdots \\
        \vec{x}_N^{\top}
    \end{pmatrix}
    \in
    \mathbb{R}^{N\times D}.
\]

Der Zielvektor ist:

\[
    \vec{t}
    =
    \begin{pmatrix}
        t_1 \\
        t_2 \\
        \vdots \\
        t_N
    \end{pmatrix}.
\]

Die Aktivierungen sind:

\[
    \vec{a}
    =
    X\vec{w}.
\]

Die Vorhersagen sind:

\[
    \vec{y}
    =
    \sigma(\vec{a}),
\]

wobei \(\sigma\) komponentenweise angewendet wird:

\[
    y_n
    =
    \sigma(a_n).
\]

Dann lautet der Gradient:

\[
    \nabla_{\vec{w}}E(\vec{w})
    =
    X^{\top}(\vec{y}-\vec{t}).
\]

\begin{definitionbox}
In Matrixform ist der Gradient:

\[
    \nabla_{\vec{w}}E(\vec{w})
    =
    X^{\top}(\vec{y}-\vec{t}).
\]
\end{definitionbox}

\begin{notebox}
Die Dimensionen sind:

\[
    X \in \mathbb{R}^{N\times D},
    \qquad
    \vec{y}-\vec{t} \in \mathbb{R}^{N},
    \qquad
    X^{\top}(\vec{y}-\vec{t}) \in \mathbb{R}^{D}.
\]
\end{notebox}

\section{Warum gibt es keine geschlossene Lösung?}

Bei linearer Regression mit quadratischem Fehler konnten wir die Normalengleichung herleiten:

\[
    \Phi^{\top}\Phi\vec{w}
    =
    \Phi^{\top}\vec{t}.
\]

Daraus folgte eine geschlossene Lösung.

Bei logistischer Regression ist die Situation anders.
Die Fehlerfunktion enthält die Sigmoid-Funktion und Logarithmen:

\[
    E(\vec{w})
    =
    -\sum_{n=1}^{N}
    \left[
        t_n\log \sigma(\vec{w}^{\top}\vec{x}_n)
        +
        (1-t_n)\log(1-\sigma(\vec{w}^{\top}\vec{x}_n))
    \right].
\]

Der Gradient ist zwar einfach:

\[
    \nabla_{\vec{w}}E(\vec{w})
    =
    X^{\top}(\vec{y}-\vec{t}),
\]

aber \(\vec{y}\) hängt nichtlinear von \(\vec{w}\) ab:

\[
    \vec{y}
    =
    \sigma(X\vec{w}).
\]

Die Gleichung

\[
    X^{\top}(\vec{y}-\vec{t})
    =
    \vec{0}
\]

ist daher nichtlinear in \(\vec{w}\).

\begin{conceptbox}
Logistische Regression hat im Allgemeinen keine geschlossene Lösung wie lineare Regression.

Man verwendet iterative Optimierungsverfahren, zum Beispiel Gradientenabstieg (gradient descent), Newton-Verfahren oder quasi-Newton-Verfahren.
\end{conceptbox}

\section{Gradientenabstieg}

Gradientenabstieg (gradient descent) aktualisiert die Gewichte in Richtung des negativen Gradienten:

\[
    \vec{w}^{(\tau+1)}
    =
    \vec{w}^{(\tau)}
    -
    \eta
    \nabla_{\vec{w}}E(\vec{w}^{(\tau)}).
\]

Dabei ist

\[
    \eta>0
\]

die Lernrate (learning rate).

Für logistische Regression:

\[
    \vec{w}^{(\tau+1)}
    =
    \vec{w}^{(\tau)}
    -
    \eta
    X^{\top}
    (
        \vec{y}^{(\tau)}-\vec{t}
    ).
\]

Dabei ist:

\[
    \vec{y}^{(\tau)}
    =
    \sigma(X\vec{w}^{(\tau)}).
\]

\begin{definitionbox}
Batch-Gradientenabstieg für logistische Regression:

\[
    \vec{w}^{(\tau+1)}
    =
    \vec{w}^{(\tau)}
    -
    \eta
    X^{\top}
    (
        \sigma(X\vec{w}^{(\tau)})-\vec{t}
    ).
\]
\end{definitionbox}

\section{Stochastischer Gradientenabstieg}

Beim stochastischen Gradientenabstieg (stochastic gradient descent, SGD) wird nicht der ganze Datensatz verwendet.
Stattdessen wird ein einzelner Datenpunkt oder ein Mini-Batch verwendet.

Für einen einzelnen Datenpunkt lautet das Update:

\[
    \vec{w}^{(\tau+1)}
    =
    \vec{w}^{(\tau)}
    -
    \eta
    (y_n-t_n)\vec{x}_n.
\]

Dabei ist:

\[
    y_n
    =
    \sigma((\vec{w}^{(\tau)})^{\top}\vec{x}_n).
\]

\begin{conceptbox}
SGD ist wichtig, weil es bei großen Datensätzen effizienter ist als Batch-Gradientenabstieg.

Viele moderne Modelle werden mit Mini-Batch-Varianten von SGD trainiert.
\end{conceptbox}

\section{Hesse-Matrix und Konvexität}

Die Fehlerfunktion der logistischen Regression ist konvex.
Das bedeutet, dass jedes lokale Minimum auch ein globales Minimum ist.

Die Hesse-Matrix (Hessian matrix) ist:

\[
    H
    =
    \nabla_{\vec{w}}^2E(\vec{w}).
\]

Für logistische Regression gilt:

\[
    H
    =
    X^{\top}RX,
\]

wobei \(R\) eine Diagonalmatrix ist:

\[
    R
    =
    \mathrm{diag}
    \left(
        y_1(1-y_1),
        y_2(1-y_2),
        \dots,
        y_N(1-y_N)
    \right).
\]

Da

\[
    0 < y_n(1-y_n) \leq \frac{1}{4},
\]

ist \(R\) positiv semidefinit.
Damit ist auch

\[
    X^{\top}RX
\]

positiv semidefinit.

\begin{definitionbox}
Die Hesse-Matrix der logistischen Regression ist:

\[
    H
    =
    X^{\top}RX,
\]

mit

\[
    R_{nn}
    =
    y_n(1-y_n).
\]
\end{definitionbox}

\begin{conceptbox}
Die Kreuzentropie der logistischen Regression ist konvex in \(\vec{w}\).

Das bedeutet:
Es gibt keine schlechten lokalen Minima.
Trotzdem muss das Minimum numerisch gefunden werden.
\end{conceptbox}

\section{Newton-Verfahren und IRLS}

Neben Gradientenabstieg kann man auch das Newton-Verfahren (Newton's method) verwenden.
Das Update lautet:

\[
    \vec{w}^{(\tau+1)}
    =
    \vec{w}^{(\tau)}
    -
    H^{-1}
    \nabla_{\vec{w}}E(\vec{w}^{(\tau)}).
\]

Für logistische Regression:

\[
    \vec{w}^{(\tau+1)}
    =
    \vec{w}^{(\tau)}
    -
    (X^{\top}RX)^{-1}
    X^{\top}(\vec{y}-\vec{t}).
\]

Dieses Verfahren kann auch als iterativ gewichtete kleinste Quadrate (iteratively reweighted least squares, IRLS) interpretiert werden.

\begin{notebox}
Newton-Verfahren kann schneller konvergieren als Gradientenabstieg, ist aber pro Iteration teurer, weil eine Matrix invertiert oder ein lineares Gleichungssystem gelöst werden muss.
\end{notebox}

\section{Entscheidungsgrenze}

Die Entscheidung des Modells basiert auf:

\[
    y(\vec{x},\vec{w})
    =
    \sigma(\vec{w}^{\top}\vec{x}).
\]

Mit Schwelle \(0{,}5\):

\[
    \hat{t}=1
    \quad
    \Longleftrightarrow
    \quad
    y(\vec{x},\vec{w})\geq 0{,}5.
\]

Da

\[
    \sigma(a)\geq 0{,}5
    \quad
    \Longleftrightarrow
    \quad
    a\geq 0,
\]

gilt:

\[
    \hat{t}=1
    \quad
    \Longleftrightarrow
    \quad
    \vec{w}^{\top}\vec{x}\geq 0.
\]

Die Entscheidungsgrenze ist daher:

\[
    \vec{w}^{\top}\vec{x}=0.
\]

\begin{definitionbox}
Die Entscheidungsgrenze (decision boundary) der logistischen Regression ist eine Hyperebene:

\[
    \vec{w}^{\top}\vec{x}=0.
\]
\end{definitionbox}

Für zwei Eingabedimensionen mit Bias-Term:

\[
    \vec{x}
    =
    \begin{pmatrix}
        1 \\
        x_1 \\
        x_2
    \end{pmatrix},
\]

ist die Entscheidungsgrenze:

\[
    w_0+w_1x_1+w_2x_2=0.
\]

Falls

\[
    w_2 \neq 0,
\]

kann man nach \(x_2\) auflösen:

\[
    x_2
    =
    -\frac{w_0}{w_2}
    -
    \frac{w_1}{w_2}x_1.
\]

Das ist eine Gerade.

\begin{conceptbox}
Logistische Regression ist ein linearer Klassifikator.

Die Wahrscheinlichkeiten sind nichtlinear, aber die Entscheidungsgrenze ist linear im Merkmalsraum.
\end{conceptbox}

\section{Nichtlineare Entscheidungsgrenzen mit Basisfunktionen}

Wie bei linearer Regression kann man Basisfunktionen verwenden.
Statt direkt \(\vec{x}\) zu verwenden, nutzt man

\[
    \vec{\phi}(\vec{x}).
\]

Das Modell lautet:

\[
    y(\vec{x},\vec{w})
    =
    \sigma(\vec{w}^{\top}\vec{\phi}(\vec{x})).
\]

Die Entscheidungsgrenze ist:

\[
    \vec{w}^{\top}\vec{\phi}(\vec{x})=0.
\]

Im transformierten Merkmalsraum ist die Grenze linear.
Im ursprünglichen Eingaberaum kann sie nichtlinear sein.

\begin{examplebox}
Für eine eindimensionale Eingabe kann man polynomiale Basisfunktionen verwenden:

\[
    \vec{\phi}(x)
    =
    \begin{pmatrix}
        1 \\
        x \\
        x^2
    \end{pmatrix}.
\]

Dann ist:

\[
    y(x,\vec{w})
    =
    \sigma(w_0+w_1x+w_2x^2).
\]

Die Entscheidungsgrenze ist:

\[
    w_0+w_1x+w_2x^2=0.
\]

Diese Gleichung kann bis zu zwei Lösungen haben.
\end{examplebox}

\section{Regularisierte logistische Regression}

Wie bei linearer Regression kann logistische Regression überanpassen.
Daher verwendet man Regularisierung.

Mit \(L_2\)-Regularisierung lautet die Fehlerfunktion:

\[
    E_{\mathrm{reg}}(\vec{w})
    =
    E(\vec{w})
    +
    \frac{\lambda}{2}
    \vec{w}^{\top}\vec{w}.
\]

Dabei ist:

\[
    E(\vec{w})
    =
    -\sum_{n=1}^{N}
    \left[
        t_n\log y_n
        +
        (1-t_n)\log(1-y_n)
    \right].
\]

\begin{definitionbox}
Regularisierte logistische Regression mit \(L_2\)-Regularisierung minimiert:

\[
    E_{\mathrm{reg}}(\vec{w})
    =
    -\sum_{n=1}^{N}
    \left[
        t_n\log y_n
        +
        (1-t_n)\log(1-y_n)
    \right]
    +
    \frac{\lambda}{2}
    \vec{w}^{\top}\vec{w}.
\]
\end{definitionbox}

Der Gradient ist:

\[
    \nabla_{\vec{w}}E_{\mathrm{reg}}(\vec{w})
    =
    X^{\top}(\vec{y}-\vec{t})
    +
    \lambda\vec{w}.
\]

\begin{notebox}
Oft wird der Bias-Term \(w_0\) nicht regularisiert.
Dann verwendet man statt \(\vec{w}^{\top}\vec{w}\) nur die Summe über die Gewichte ohne Bias.
\end{notebox}

\section{MAP-Interpretation der Regularisierung}

Regularisierung kann probabilistisch als Prior interpretiert werden.
Wir verwenden einen Gauß-Prior auf die Gewichte:

\[
    p(\vec{w})
    =
    \mathcal{N}
    \left(
        \vec{w}
        \mid
        \vec{0},
        \alpha^{-1}I
    \right).
\]

Dieser Prior hat die Form:

\[
    p(\vec{w})
    \propto
    \exp
    \left(
        -\frac{\alpha}{2}
        \vec{w}^{\top}\vec{w}
    \right).
\]

Der negative Log-Prior ist bis auf Konstanten:

\[
    -\log p(\vec{w})
    =
    \frac{\alpha}{2}
    \vec{w}^{\top}\vec{w}
    +
    \text{const}.
\]

Die MAP-Schätzung minimiert:

\[
    -\log p(\vec{t}\mid X,\vec{w})
    -
    \log p(\vec{w}).
\]

Das ergibt:

\[
    E_{\mathrm{MAP}}(\vec{w})
    =
    E(\vec{w})
    +
    \frac{\alpha}{2}
    \vec{w}^{\top}\vec{w}
    +
    \text{const}.
\]

\begin{conceptbox}
\(L_2\)-regularisierte logistische Regression ist MAP-Schätzung mit Bernoulli-Likelihood und Gauß-Prior auf die Gewichte.

\[
    \text{Bernoulli-Likelihood}
    +
    \text{Gauß-Prior}
    \quad
    \Longrightarrow
    \quad
    \text{Kreuzentropie}
    +
    L_2\text{-Regularisierung}.
\]
\end{conceptbox}

\section{Numerische Stabilität}

In Implementierungen muss man auf numerische Stabilität achten.
Die Sigmoid-Funktion

\[
    \sigma(a)
    =
    \frac{1}{1+\exp(-a)}
\]

kann für sehr große negative \(a\) numerisch problematisch sein, weil

\[
    \exp(-a)
\]

sehr groß wird.

Auch die Kreuzentropie

\[
    -\left[
        t\log y
        +
        (1-t)\log(1-y)
    \right]
\]

ist problematisch, wenn \(y\) exakt \(0\) oder \(1\) wird, weil

\[
    \log(0)
\]

nicht definiert ist.

\begin{notebox}
In praktischen Implementierungen verwendet man oft numerisch stabile Varianten der Sigmoid- und Kreuzentropie-Berechnung.

Außerdem werden Wahrscheinlichkeiten manchmal auf ein kleines Intervall beschränkt, zum Beispiel

\[
    y \in [\varepsilon,1-\varepsilon].
\]
\end{notebox}

\section{Evaluation eines binären Klassifikators}

Nach dem Training muss das Modell bewertet werden.
Dazu verwendet man Testdaten:

\[
    \mathcal{D}_{\mathrm{test}}
    =
    \{(\vec{x}_1^{\,\mathrm{test}},t_1^{\mathrm{test}}),\dots,(\vec{x}_M^{\,\mathrm{test}},t_M^{\mathrm{test}})\}.
\]

Für jeden Testpunkt berechnet man:

\[
    y_m
    =
    p(t=1\mid\vec{x}_m^{\,\mathrm{test}},\vec{w}).
\]

Dann wählt man eine Schwelle, meistens:

\[
    \tau = 0{,}5.
\]

Die vorhergesagte Klasse ist:

\[
    \hat{t}_m
    =
    \begin{cases}
        1, & y_m \geq \tau, \\
        0, & y_m < \tau.
    \end{cases}
\]

\section{Accuracy}

Die Accuracy (accuracy) ist der Anteil korrekt klassifizierter Datenpunkte.

\[
    \mathrm{accuracy}
    =
    \frac{
        \text{Anzahl korrekter Vorhersagen}
    }{
        \text{Anzahl aller Vorhersagen}
    }.
\]

Mathematisch:

\[
    \mathrm{accuracy}
    =
    \frac{1}{M}
    \sum_{m=1}^{M}
    \mathbb{1}(\hat{t}_m=t_m).
\]

\begin{definitionbox}
Die Accuracy ist:

\[
    \mathrm{accuracy}
    =
    \frac{1}{M}
    \sum_{m=1}^{M}
    \mathbb{1}(\hat{t}_m=t_m).
\]

Dabei ist \(\mathbb{1}(\hat{t}_m=t_m)\) gleich \(1\), wenn die Vorhersage korrekt ist, und sonst \(0\).
\end{definitionbox}

\begin{notebox}
Accuracy kann irreführend sein, wenn die Klassen stark unausgeglichen sind.

Wenn \(95\%\) der Daten Klasse \(0\) sind, erreicht ein Modell, das immer \(0\) vorhersagt, bereits \(95\%\) Accuracy.
\end{notebox}

\section{Konfusionsmatrix}

Für binäre Klassifikation unterscheidet man:

\begin{itemize}
    \item True Positive (TP): \(t=1\) und \(\hat{t}=1\),
    \item True Negative (TN): \(t=0\) und \(\hat{t}=0\),
    \item False Positive (FP): \(t=0\), aber \(\hat{t}=1\),
    \item False Negative (FN): \(t=1\), aber \(\hat{t}=0\).
\end{itemize}

Die Konfusionsmatrix (confusion matrix) ist:

\[
\begin{array}{c|cc}
    & \hat{t}=1 & \hat{t}=0 \\
    \hline
    t=1 & \mathrm{TP} & \mathrm{FN} \\
    t=0 & \mathrm{FP} & \mathrm{TN}
\end{array}
\]

\begin{definitionbox}
Die Konfusionsmatrix zeigt, welche Klassen korrekt oder falsch vorhergesagt wurden.
Sie ist besonders hilfreich, wenn Accuracy allein nicht aussagekräftig ist.
\end{definitionbox}

\section{Precision, Recall und F1-Score}

Aus der Konfusionsmatrix lassen sich weitere Maße berechnen.

Precision (precision):

\[
    \mathrm{precision}
    =
    \frac{\mathrm{TP}}{\mathrm{TP}+\mathrm{FP}}.
\]

Recall (recall):

\[
    \mathrm{recall}
    =
    \frac{\mathrm{TP}}{\mathrm{TP}+\mathrm{FN}}.
\]

F1-Score:

\[
    F_1
    =
    2
    \frac{
        \mathrm{precision}\cdot\mathrm{recall}
    }{
        \mathrm{precision}+\mathrm{recall}
    }.
\]

\begin{conceptbox}
Precision fragt:

\[
    \text{Wie viele positive Vorhersagen waren wirklich positiv?}
\]

Recall fragt:

\[
    \text{Wie viele wirklich positive Fälle wurden gefunden?}
\]
\end{conceptbox}

\section{Wahl der Entscheidungsschwelle}

Die Standardschwelle ist:

\[
    \tau=0{,}5.
\]

Manchmal ist eine andere Schwelle sinnvoll.
Die allgemeine Entscheidungsregel lautet:

\[
    \hat{t}
    =
    \begin{cases}
        1, & y \geq \tau, \\
        0, & y < \tau.
    \end{cases}
\]

Wenn False Positives teuer sind, wählt man oft eine höhere Schwelle.
Wenn False Negatives teuer sind, wählt man oft eine niedrigere Schwelle.

\begin{examplebox}
Bei einer medizinischen Voruntersuchung kann es schlimmer sein, eine kranke Person fälschlich als gesund einzustufen, als eine gesunde Person zur weiteren Untersuchung zu schicken.

Dann möchte man False Negatives reduzieren und kann eine niedrigere Schwelle \(\tau\) wählen.
\end{examplebox}

\section{Mehrklassen-Klassifikation}

Logistische Regression kann auf mehr als zwei Klassen erweitert werden.
Für \(K\) Klassen gilt:

\[
    t \in \{1,2,\dots,K\}.
\]

Für jede Klasse \(k\) berechnet man einen Score:

\[
    a_k
    =
    \vec{w}_k^{\top}\vec{x}.
\]

Diese Scores werden mit der Softmax-Funktion in Wahrscheinlichkeiten umgewandelt.

\begin{definitionbox}
Die Softmax-Funktion (softmax function) ist:

\[
    y_k
    =
    p(t=k\mid\vec{x})
    =
    \frac{
        \exp(a_k)
    }{
        \sum_{j=1}^{K}
        \exp(a_j)
    }.
\]
\end{definitionbox}

Die Wahrscheinlichkeiten erfüllen:

\[
    y_k \geq 0
\]

und

\[
    \sum_{k=1}^{K} y_k = 1.
\]

\begin{conceptbox}
Die Softmax-Funktion ist die Mehrklassen-Verallgemeinerung der Sigmoid-Funktion.

Sie wandelt mehrere Scores in eine Wahrscheinlichkeitsverteilung über Klassen um.
\end{conceptbox}

\section{Mehrklassen-Kreuzentropie}

Für Mehrklassen-Klassifikation verwendet man oft One-Hot-Kodierung (one-hot encoding).
Der Zielvektor ist:

\[
    \vec{t}
    =
    \begin{pmatrix}
        t_1 \\
        \vdots \\
        t_K
    \end{pmatrix},
\]

wobei genau eine Komponente \(1\) ist und alle anderen \(0\) sind.

Für einen Datenpunkt ist die Kreuzentropie:

\[
    E_n
    =
    -\sum_{k=1}^{K}
    t_{nk}\log y_{nk}.
\]

Für den ganzen Datensatz:

\[
    E
    =
    -\sum_{n=1}^{N}
    \sum_{k=1}^{K}
    t_{nk}\log y_{nk}.
\]

\begin{definitionbox}
Die Mehrklassen-Kreuzentropie (multiclass cross-entropy) ist:

\[
    E
    =
    -\sum_{n=1}^{N}
    \sum_{k=1}^{K}
    t_{nk}\log y_{nk}.
\]
\end{definitionbox}

\section{Softmax als probabilistisches Modell}

Die Softmax-Ausgabe wird interpretiert als:

\[
    y_k
    =
    p(t=k\mid\vec{x}).
\]

Die vorhergesagte Klasse ist:

\[
    \hat{t}
    =
    \arg\max_{k}
    y_k.
\]

Da Softmax monoton in den Scores ist, gilt auch:

\[
    \hat{t}
    =
    \arg\max_{k}
    a_k.
\]

\begin{notebox}
Für die Entscheidung reicht der größte Score.
Für Wahrscheinlichkeiten und Training braucht man aber die Softmax-normalisierten Werte.
\end{notebox}

\section{Zusammenhang zu neuronalen Netzen}

Logistische Regression kann als einfachstes neuronales Netz betrachtet werden:

\[
    \vec{x}
    \longrightarrow
    \vec{w}^{\top}\vec{x}
    \longrightarrow
    \sigma(\vec{w}^{\top}\vec{x}).
\]

Es gibt keine versteckte Schicht (hidden layer).
Das Modell ist daher ein linearer Klassifikator mit nichtlinearer Ausgabefunktion.

\begin{conceptbox}
Ein neuronales Netz für binäre Klassifikation verwendet am Ende oft genau dieselbe Kombination:

\[
    \text{Logit}
    \longrightarrow
    \text{Sigmoid}
    \longrightarrow
    \text{Kreuzentropie}.
\]

Logistische Regression ist daher ein wichtiger Spezialfall neuronaler Netze.
\end{conceptbox}

\section{Praktisches Beispiel: Binäre Zielvariable erzeugen}

In vielen Datensätzen ist der ursprüngliche Zielwert kontinuierlich.
Für binäre Klassifikation kann man daraus eine binäre Zielvariable erzeugen.

Angenommen, ein Zielwert \(z_n\) ist kontinuierlich.
Wir wählen eine Schwelle \(c\) und definieren:

\[
    t_n
    =
    \begin{cases}
        1, & z_n \geq c, \\
        0, & z_n < c.
    \end{cases}
\]

\begin{examplebox}
Wenn \(z_n\) ein Wohnungspreis ist, kann man definieren:

\[
    t_n
    =
    \begin{cases}
        1, & \text{Preis ist hoch}, \\
        0, & \text{Preis ist niedrig}.
    \end{cases}
\]

Dann kann logistische Regression verwendet werden, um die Wahrscheinlichkeit für einen hohen Preis zu modellieren.
\end{examplebox}

\begin{notebox}
Das Umwandeln eines kontinuierlichen Zielwerts in eine binäre Klasse kann nützlich sein, verliert aber Information.

Ein Preis von \(101\) und ein Preis von \(1000\) können nach der Schwelle beide dieselbe Klasse bekommen.
\end{notebox}

\section{Feature-Skalierung}

Wie bei linearer Regression ist Feature-Skalierung wichtig.
Wenn Merkmale sehr unterschiedliche Größenordnungen haben, kann Gradientenabstieg schlecht funktionieren.

Eine Standardisierung ist:

\[
    x_j'
    =
    \frac{x_j-\mu_j}{\sigma_j}.
\]

Dabei sind \(\mu_j\) und \(\sigma_j\) Mittelwert und Standardabweichung des Merkmals im Trainingsdatensatz.

\begin{examtipbox}
Wichtig:
Skalierungsparameter werden nur aus Trainingsdaten berechnet.

Dieselben Parameter werden dann auf Validierungs- und Testdaten angewendet.
\end{examtipbox}

\section{Häufige Fehler}

\begin{examtipbox}
Typische Fehler bei logistischer Regression:

\begin{enumerate}
    \item Denken, logistische Regression sei ein Regressionsmodell für kontinuierliche Zielwerte.
    \item Die Ausgabe \(\vec{w}^{\top}\vec{x}\) direkt als Wahrscheinlichkeit interpretieren.
    \item Vergessen, die Sigmoid-Funktion anzuwenden.
    \item \(p(t=1\mid\vec{x})\) und \(p(t=0\mid\vec{x})\) verwechseln.
    \item Kreuzentropie mit quadratischem Fehler ersetzen, ohne die probabilistische Bedeutung zu beachten.
    \item Die Ableitung der Sigmoid-Funktion falsch berechnen.
    \item Beim Gradienten das einfache Ergebnis \((y_n-t_n)\vec{x}_n\) nicht erkennen.
    \item Die Entscheidungsgrenze mit der Sigmoid-Kurve verwechseln.
    \item Accuracy verwenden, obwohl die Klassen stark unausgeglichen sind.
    \item Bei MAP-Regularisierung den Prior vergessen.
    \item Den Bias-Term versehentlich regularisieren, obwohl dies nicht gewünscht ist.
\end{enumerate}
\end{examtipbox}

\section{Mini-Aufgaben}

\subsection{Aufgabe 1: Sigmoid berechnen}

Berechne

\[
    \sigma(0).
\]

\begin{examplebox}
Mit

\[
    \sigma(a)
    =
    \frac{1}{1+\exp(-a)}
\]

folgt:

\[
    \sigma(0)
    =
    \frac{1}{1+\exp(0)}
    =
    \frac{1}{1+1}
    =
    \frac{1}{2}.
\]
\end{examplebox}

\subsection{Aufgabe 2: Wahrscheinlichkeit und Gegenwahrscheinlichkeit}

Ein logistisches Regressionsmodell gibt aus:

\[
    p(t=1\mid\vec{x})=0{,}8.
\]

Berechne

\[
    p(t=0\mid\vec{x}).
\]

\begin{examplebox}
Da es nur zwei Klassen gibt:

\[
    p(t=0\mid\vec{x})
    =
    1-p(t=1\mid\vec{x}).
\]

Also:

\[
    p(t=0\mid\vec{x})
    =
    1-0{,}8
    =
    0{,}2.
\]
\end{examplebox}

\subsection{Aufgabe 3: Binäre Kreuzentropie}

Für einen Datenpunkt gilt:

\[
    t=1,
    \qquad
    y=0{,}9.
\]

Berechne den Kreuzentropie-Verlust.

\begin{examplebox}
Der Verlust ist:

\[
    E
    =
    -\left[
        t\log y
        +
        (1-t)\log(1-y)
    \right].
\]

Mit \(t=1\):

\[
    E
    =
    -\log y.
\]

Also:

\[
    E
    =
    -\log(0{,}9).
\]
\end{examplebox}

\subsection{Aufgabe 4: Gradient für einen Datenpunkt}

Gegeben seien

\[
    y_n=0{,}7,
    \qquad
    t_n=1,
\]

und

\[
    \vec{x}_n
    =
    \begin{pmatrix}
        1 \\
        2
    \end{pmatrix}.
\]

Berechne den Gradienten

\[
    \nabla_{\vec{w}}E_n(\vec{w}).
\]

\begin{examplebox}
Der Gradient ist:

\[
    \nabla_{\vec{w}}E_n(\vec{w})
    =
    (y_n-t_n)\vec{x}_n.
\]

Einsetzen:

\[
    y_n-t_n
    =
    0{,}7-1
    =
    -0{,}3.
\]

Also:

\[
    \nabla_{\vec{w}}E_n(\vec{w})
    =
    -0{,}3
    \begin{pmatrix}
        1 \\
        2
    \end{pmatrix}
    =
    \begin{pmatrix}
        -0{,}3 \\
        -0{,}6
    \end{pmatrix}.
\]
\end{examplebox}

\subsection{Aufgabe 5: Entscheidungsgrenze}

Ein Modell hat

\[
    w_0=-1,
    \qquad
    w_1=2.
\]

Für eine eindimensionale Eingabe mit Bias gilt:

\[
    a
    =
    w_0+w_1x.
\]

Bestimme die Entscheidungsgrenze.

\begin{examplebox}
Die Entscheidungsgrenze ist:

\[
    a=0.
\]

Also:

\[
    -1+2x=0.
\]

Damit:

\[
    2x=1
\]

und

\[
    x=\frac{1}{2}.
\]

Für \(x\geq \frac{1}{2}\) sagt das Modell Klasse \(1\) voraus.
Für \(x<\frac{1}{2}\) sagt es Klasse \(0\) voraus.
\end{examplebox}

\section{Zusammenfassung}

\begin{itemize}
    \item Logistische Regression (logistic regression) ist ein Modell für binäre Klassifikation.
    \item Das Modell berechnet zuerst einen linearen Score:
    \[
        a
        =
        \vec{w}^{\top}\vec{x}.
    \]
    \item Die Sigmoid-Funktion wandelt den Score in eine Wahrscheinlichkeit um:
    \[
        y
        =
        \sigma(a)
        =
        \frac{1}{1+\exp(-a)}.
    \]
    \item Die Modellwahrscheinlichkeit ist:
    \[
        p(t=1\mid\vec{x},\vec{w})
        =
        \sigma(\vec{w}^{\top}\vec{x}).
    \]
    \item Für die Gegenklasse gilt:
    \[
        p(t=0\mid\vec{x},\vec{w})
        =
        1-y.
    \]
    \item Der Logit ist linear:
    \[
        \log
        \left(
            \frac{y}{1-y}
        \right)
        =
        \vec{w}^{\top}\vec{x}.
    \]
    \item Die Bernoulli-Likelihood ist:
    \[
        L(\vec{w})
        =
        \prod_{n=1}^{N}
        y_n^{t_n}
        (1-y_n)^{1-t_n}.
    \]
    \item Die negative Log-Likelihood ist die binäre Kreuzentropie:
    \[
        E(\vec{w})
        =
        -\sum_{n=1}^{N}
        \left[
            t_n\log y_n
            +
            (1-t_n)\log(1-y_n)
        \right].
    \]
    \item Die Ableitung der Sigmoid-Funktion ist:
    \[
        \sigma'(a)
        =
        \sigma(a)(1-\sigma(a)).
    \]
    \item Der Gradient für einen Datenpunkt ist:
    \[
        \nabla_{\vec{w}}E_n(\vec{w})
        =
        (y_n-t_n)\vec{x}_n.
    \]
    \item Der Gradient für den Datensatz ist:
    \[
        \nabla_{\vec{w}}E(\vec{w})
        =
        X^{\top}(\vec{y}-\vec{t}).
    \]
    \item Die Entscheidungsgrenze ist:
    \[
        \vec{w}^{\top}\vec{x}=0.
    \]
    \item \(L_2\)-regularisierte logistische Regression entspricht MAP mit Gauß-Prior.
    \item Für mehrere Klassen verwendet man die Softmax-Funktion:
    \[
        y_k
        =
        \frac{\exp(a_k)}
        {\sum_{j=1}^{K}\exp(a_j)}.
    \]
\end{itemize}

\section{Typische Prüfungsfragen}

\begin{examtipbox}
Mögliche Prüfungsfragen zu diesem Kapitel:

\begin{enumerate}
    \item Warum ist lineare Regression nicht direkt für Klassifikation geeignet?
    \item Definiere die Sigmoid-Funktion.
    \item Zeige, dass \(\sigma(0)=\frac{1}{2}\).
    \item Zeige, dass \(\sigma'(a)=\sigma(a)(1-\sigma(a))\).
    \item Wie modelliert logistische Regression \(p(t=1\mid\vec{x})\)?
    \item Was ist der Logit?
    \item Zeige, dass logistische Regression lineare Log-Odds modelliert.
    \item Schreibe die Bernoulli-Likelihood für logistische Regression auf.
    \item Leite die Log-Likelihood her.
    \item Warum ist die Kreuzentropie die negative Log-Likelihood?
    \item Leite den Gradienten für einen Datenpunkt her.
    \item Leite den Gradienten in Matrixform her.
    \item Warum gibt es keine geschlossene Lösung wie bei linearer Regression?
    \item Wie lautet das Update des Batch-Gradientenabstiegs?
    \item Was ist die Entscheidungsgrenze der logistischen Regression?
    \item Warum ist logistische Regression ein linearer Klassifikator?
    \item Wie kann man nichtlineare Entscheidungsgrenzen erzeugen?
    \item Wie funktioniert \(L_2\)-Regularisierung bei logistischer Regression?
    \item Wie hängt \(L_2\)-Regularisierung mit MAP zusammen?
    \item Was ist der Unterschied zwischen Sigmoid und Softmax?
    \item Wie lautet die Mehrklassen-Kreuzentropie?
    \item Warum kann Accuracy bei unausgeglichenen Klassen irreführend sein?
\end{enumerate}
\end{examtipbox}